\section{Conclusion}\label{sec:conclusion}
This paper demonstrates that co-designing a DBMS and an SSD can improve throughput and extend SSD lifespan.
Our four out-of-place schemes reshape the DBMS write pattern to mitigate WA at both the DB and SSD layers. % ---(1) compression, (2) GDT, (3) aligning the GC unit, and (4) the NoWA---
We further show that the design naturally leverages ZNS and FDP.
\revision[R1.W1]{Future work includes extending these techniques to LSM engines, supporting multi- or shared-device deployments, and exploring disk- and HM-SMR--based systems.}

Beyond DBMSs and SSDs, our findings apply broadly to storage systems.
Our large sequential write patterns also benefit disk-based systems~\cite{DBLP:journals/corr/abs-2205-11753,DBLP:conf/sosp/BornholtJACKMSS21,skylight}.
Systems with their own translation layer and garbage collection can adopt our DB-level WAF optimizations, 
while log-structured filesystems can use the NoWA pattern to reduce aging~\cite{DBLP:conf/hotstorage/ParkE22,DBLP:conf/usenix/KadekodiNG18}.
We hope this work helps storage designers better understand the implications of system writes for devices and how to fully exploit device capabilities.
